\section{Experiments} \label{sec:experiments}
This section contains details of all the performed experiments, as well as 
other considerations related thereto.
\subsection{Experiment Suite}
All models were trained over data from an 11 year time-frame, from January 2015 
to mid-March 2026. They were also trained in a walk-forward sliding-window with 
an initial training window of two years, then a test window of 30 days, as well 
as a step size of 30 days. As previously mentioned, data-groups were 
constructed in 
order to create different data-scenarios for training. These groups consisted 
of the following categories: 
\begin{itemize}
	\item \textbf{Weather}: All weather-related features. 
	\item \textbf{WeatherLags}: T-Minus 24 hour data of the Weather group. 
	\item \textbf{Grid}: All data related to \gls{pc}, aside from the foreign 
	energy exchange features.
	\item \textbf{GridLags}: T-Minus 24 hour data of the Grid group. 
	\item \textbf{GridExchange}: Only the data related to foreign energy 
	exchange. 
	\item \textbf{GridExchangeLags}: T-Minus 24 hour data of the GridExchange 
	group. 
	\item \textbf{Price}: Columns containing the absolute price and price 
	delta. These are not included in the training data, but are instead 
	individually used as prediction targets depending on the type of experiment 
	being run, either absolute price or price delta. 
	\item \textbf{PriceLags}: T-Minus 24, 48, and 
	168 hour data of both absolute 
	price and price delta. 
	\item \textbf{Time}: Data related to time, such as hour, day of the week, 
	day of the month, day of the year, as well as sine and cosine 
	representations.
	\item \textbf{AllFeatures}: Simply a collection of all previous groupings.
\end{itemize}

Based on these groups, we created 13 base experiments, generally increasing 
information sequentially: 
\begin{enumerate}
	\item Weather, Time
	\item Weather, WeatherLags, Time
	\item Weather, PriceLags, Time
	\item Weather, WeatherLags, PriceLags, Time
	\item Weather, Grid, GridExchange, Time
	\item Weather, WeatherLags, Grid, 
	GridExchange, Time
	\item Weather, Grid, GridExchange, PriceLags, 
	Time
	\item Weather, WeatherLags, Grid, 
	GridExchange, PriceLags, Time
	\item Weather, Grid, GridExchange, GridLags, 
	GridExchangeLags, Time
	\item Weather, WeatherLags, Grid, 
	GridExchange, GridLags,	GridExchangeLags, 
	Time
	\item Weather, Grid, GridExchange, GridLags, 
	GridExchangeLags, PriceLags, Time
	\item Weather, WeatherLags, Grid, 
	GridExchange, GridLags, GridExchangeLags, 
	PriceLags, Time
	\item AllFeatures
\end{enumerate}
These 13 experiments were run at horizons 0, 24, 48, 72, 96, 120, 144, and 168 
hour for both targets absolute price and price delta. Note that price delta was 
not used as a target for horizon 0 hour, as the delta between the current and 
target would always be 0. The 24 hour horizon was run for all eight models, 
with the subsequent horizons being run for all tree-based models, but only for 
select neural network-based models, due to their training-time constraints. 
This resulted in $4 \times 13 \times (7+8) = 780$ tree-based models, where 4 is 
the number of tree-based model types, 13 is the number of feature set 
configurations, and $(7+8)$ represents the 8 absolute price horizons and 7 
price delta horizons respectively (horizon 0h is excluded from the delta target 
as it is always zero). For the neural network-based models and AutoGluon, an 
equivalent set of 4 model types were evaluated, yielding $4 \times (13 \times 2 
+ 3 \times 7 \times 2) = 272$ models, where the first term covers all 13 
feature sets at the 24h horizon for both targets, the second term covers the 3 
selected feature sets across the remaining 7 horizons and both targets, and the 
factor of 2 in both terms represents the two prediction targets. The reduced 
number of neural network-based experiments was determined by their 
computational constraints, with the 3 selected feature sets identified as the 
best, mean, and worst performers based on their 24h horizon results. This 
results in a total of $780 + 272 = 1052$ models trained.

\subsection{Evaluation Criteria}\label{sec:evaluationCriteria}
While there are many different criteria one could chose to focus on when 
comparing the performance of models, the one chosen for this paper is 
\gls{mae}. The equation for calculating this metric is shown in 
\autoref{eq:mae}.
\begin{equation}
	\text{MAE} = \frac{1}{n} \sum_{t=1}^{n} |y_t - \hat{y}_t|
	\label{eq:mae}
\end{equation}
where $y_t$ and $\hat{y}_t$ are the ground truth target value at 
prediction time 
and the predicted value at prediction time respectively.\\ 
In this context, this metric represents the average amount of EUR/MWh 
the model was off. This metric was partly chosen due to the fact that it 
is comparable across various model architectures and target types. 
Furthermore, it is expressed in the same unit as the target, making it 
easily interpretable. Lastly, in a market with frequent price spikes, 
both positive and negative, it does not penalize these extremes a 
harshly as other criteria. Other considered metrics are discussed in 
\cref{app:A:equations}, and all of the values, both for the chosen and 
non-chosen metrics are available online in the GitHub repository. 
Additionally, a naive baseline, that simply reuses yesterday's price, 
calculated as follows: $\hat{y}_t = y_{t-24}$ was computed across all 
horizons to establish a non-ML reference point. 

\subsection{Experimental Setup}
All model training, validation, and testing were conducted locally on a machine 
running Microsoft Windows 11 Home, equipped with a 13th Gen Intel Core 
i7-13650HX processor and 16 GB of RAM, utilizing CPU computation exclusively. 
The exclusion of the GPU was evaluated on the basis that the libraries used for 
training and validation of the neural network-based models were optimised for 
the CPU. This means that while training time was a somewhat limiting factor for 
the neural network-based models in this paper, this aspect may be mitigated 
with different training approaches and hardware, however, the performance 
metrics, aside from training time, should be unaffected.

\subsection{Training, Testing \& Validation}\label{sec:split_protocol}
Model evaluation in this paper is performed using walk-forward 
validation with an expanding training window. This approach was 
chosen over standard k-fold cross-validation, which randomly 
partitions data into folds without regard for temporal ordering. 
Applying k-fold to time-series data would allow a model to be 
trained on future observations and tested on past ones, 
constituting a form of data leakage that would produce 
optimistically biased results. Walk-forward validation avoids 
this by strictly preserving the temporal ordering of time-series 
data. The model is always trained exclusively on historical 
data and evaluated on data it has never seen.

The procedure works as follows: An initial training window 
covering January 2015 through December 2016 is used to train 
each model for the first time. The model is then evaluated on 
the immediately following 30-day test window. After evaluation, 
the training window expands to include those 30 days, and the 
test window slides forward to the next 30-day period. This 
process repeats until the full dataset is exhausted. The 
expanding window design ensures that all available historical 
data is utilised at each step, while the fixed 30-day test 
window provides a consistent evaluation unit across all 
iterations. A figure illustrating this process is shown in 
Figure X. While this approach of essentially retraining the model 12 times per 
year, for every year in the training data, is computationally expensive, it is 
the generally preferred method of training when dealing with time-series 
data\hltodo{.}{needs source}. 

\subsection{Data Leakage Prevention}
The perceptive reader may have noticed that the master matrix dataset retains 
the raw spot price and all target columns to serve as a shared foundation for 
all experiments. However, feeding this matrix directly to a model would 
constitute severe data leakage, as the model could trivially learn to copy the 
target value rather than forecast it. To prevent this, a Target Shield is 
applied at load time before each training run, which drops all leaky columns 
from the feature set. Specifically, the active prediction target is removed 
from the features, the raw spot price is excluded except when included as a 
lagged historical feature, and the concurrent price delta columns are removed. 

\subsection{Feature Importance}
After training various constellations and horizons of many of the models, 
feature importance was performed on the best performing models of each 
architecture of each horizon. Feature importance was determined by performing 
permutation importance, a method of systematically dropping each column 
individually and measuring the impact on your evaluation criteria. For this 
paper, it was deemed that if the dropping of a column resulting in even a drop 
of 0,01 EUR/MWh, then it would be allowed to stay. This method is an 
extensively rigorous, though computationally expensive, approach to feature 
importance. The results of these pruned models will be discussed in 
\autoref{sec:results}.

\subsection{Supplementary Experiments}
To address generalisability and the limitations that come with running 
the models on their default settings i.e. with no hyperparameter tuning 
and with default loss functions, five supplementary experiments will be 
conducted:
\paragraph{Loss Function Alignment}: The best-performing neural network 
model, as will be identified from 24h baseline results, will be 
retrained using \gls{mae} loss instead of \gls{mse} to evaluate whether 
loss-metric alignment would improve performance.
\paragraph{Hyperparameter Tuning}: Optuna-based hyperparameter 
optimization will be applied to tree models on the best, mean, and worst 
performing feature sets at 24h for both targets, with 30 trials per 
configuration, to gauge the improvement such tuning may yield.
\paragraph{AutoGluon Quality Escalation}: AutoGluon will be allowed to 
train using the medium\_quality preset and a 180 second time limit. This 
will be done on the best feature set at 24h, and is to be compared 
against the baseline fastest preset, with a six second time limit, that 
was used on the ordinary experiments.
\paragraph{Midas Energy Data}: Internally used weather features from 
Midas Energy will be merged with the master matrix and tested on all 
models at the 24h-horizon and on the tree-based models at all horizons, 
to compare against the raw DMI station-level aggregation.
\paragraph{DK2 Generalization Check}: The best tree and best neural 
network models will be trained on DK2 data at a 24h-horizon across best, 
mean, and worst feature sets to verify that DK1 findings generalize 
regionally.\\
The results for all of these supplementary experiments may be found in 
section \hltodo{reference}{.}
